%% ============================================================
%%  PAGINA DE TITLU, COPYRIGHT, CUPRINS, CUVÂNT ÎNAINTE
%%  (comun tuturor volumelor; textul variabil vine din macro-uri)
%% ============================================================
\thispagestyle{empty}
\vspace*{2.2cm}
\begin{center}
  {\color{ink}\rule{\linewidth}{3pt}}\vspace{6mm}\\
  {\fontsize{30}{36}\selectfont\bfseries\color{ink}\titluserie}\\[4mm]
  {\color{onm}\rule{0.4\linewidth}{1.5pt}}\\[6mm]
  {\LARGE \titluvolum}\\[3mm]
  {\large\itshape\textcolor{grey}{\subtitluvolum}}\\[8mm]
  {\normalsize\textcolor{grey}{Teoria completă și probleme}}\\[10mm]
  {\color{ink}\rule{\linewidth}{3pt}}\\[8mm]
  {\large\itshape Vlad-Ștefan Constantinescu}
\end{center}
\vfill
\clearpage

%% ---------------- pagina de copyright ----------------
\thispagestyle{empty}
\vspace*{1cm}
\noindent{\small
\titluserie, \titluvolum\\
Volumul \numarvolum{} dintr-o serie de patru volume\\[2mm]
Autor: Vlad-Ștefan Constantinescu\\
Ediția a doua, București, 2026\\[3mm]
\copyright{} 2026 Vlad-Ștefan Constantinescu. Toate drepturile rezervate.\\[2mm]
Nicio parte a acestei lucrări nu poate fi reprodusă sau transmisă sub nicio formă și prin niciun mijloc,
electronic sau mecanic, fără acordul prealabil scris al autorului, cu excepția unor scurte citate
în recenzii și în lucrări de referință.\\[3mm]
Tehnoredactat de autor în \LaTeX{} (Xe\LaTeX).\\[3mm]
\ifx\isbnvolum\empty\else ISBN: \isbnvolum\fi

}
\vfill
\clearpage

\tableofcontents
\clearpage

%% ============================================================
%%  CUVÂNT ÎNAINTE
%% ============================================================
\thispagestyle{empty}
\addcontentsline{toc}{chapter}{Cuvânt înainte}
\vspace*{1.5cm}
{\noindent\color{onm}\rule{\linewidth}{3pt}}\par\vspace{3mm}
{\noindent\fontsize{28}{32}\selectfont\bfseries\color{ink} Cuvânt înainte}\par\vspace{3mm}
{\noindent\color{onm}\rule{\linewidth}{3pt}}\par\vspace{8mm}

\noindent \textbf{Despre acest volum.} Acesta este volumul \numarvolum{} dintr-o serie de patru volume.
\despreacestvolum{} Volumele care îl însoțesc sunt dedicate \celelaltevolume.

\medskip
\noindent Volumul are două mari componente: \textbf{teoria}, cu demonstrații complete și exemple, și \textbf{problemele}, grupate pe trei niveluri de dificultate (OLM $\cdot$ OJM $\cdot$ ONM). În partea de teorie am inclus două secțiuni pe care le consider esențiale și care lipsesc adesea din cărțile de olimpiadă: una de \textbf{redactare} (cum se scrie o demonstrație riguroasă și ce greșeli de formă penalizează juriul) și una de \textbf{construcții} (ce obiecte să încerci când enunțul nu îți dă o direcție).

\medskip
\noindent \textbf{Cele trei niveluri.} Pe tot parcursul cărții, problemele sunt clasificate după cele trei
etape ale Olimpiadei Naționale de Matematică: \textbf{OLM} (olimpiada locală, adică etapa pe școală și pe oraș),
\textbf{OJM} (olimpiada județeană) și \textbf{ONM} (olimpiada națională, etapa finală). În linii mari,
problemele OLM sunt probleme de concurs accesibile, OJM corespunde unei olimpiade regionale solide, iar
problemele ONM sunt de nivelul unei finale naționale: comparabile, pentru materia acestei cărți, cu cele mai
grele probleme din concursurile studențești precum Putnam.

\medskip
\noindent \textbf{Analiză vs.\ algebră.} Analiza matematică este un subiect continuu: materia clasei a XII-a nu este altceva decât o extensie firească a celei din clasa a XI-a, iar limitele, integralele și seriile trăiesc în aceeași lume. De aceea am construit o punte explicită între cele două clase. Algebra stă altfel: structurile algebrice din clasa a XII-a (grupuri, inele, corpuri, polinoame) formează un bloc coerent în sine, cu fire clare între ele, dar nu au nicio legătură cu algebra liniară din clasa a XI-a. Cele două jumătăți de algebră sunt lumi separate, și am ales să le tratez ca atare.

\medskip
\noindent \textbf{Despre metode.} Am îmbinat, deliberat, două perspective. Pe de o parte, tehnicile care se văd la olimpiadele din România, de la faza locală până la cea națională. Pe de altă parte, unele metode pe care le-am întâlnit în primii doi ani de facultate din sistemul \emph{CPGE} (\textit{classes préparatoires aux grandes écoles}) din Franța (funcții generatoare, comparări serie-integrală, asimptotică prin singularități) și pe care le-am găsit perfect utilizabile și la nivel de olimpiadă, adesea mai clare și mai puternice decât abordările ad-hoc. Puntea dintre clasele a XI-a și a XII-a este locul unde aceste două perspective se întâlnesc cel mai firesc.

\medskip
\noindent \textbf{Despre programă.} Nivelurile OLM $\cdot$ OJM $\cdot$ ONM indică \emph{dificultatea}, nu strict
programa etapei corespunzătoare. Am inclus deliberat, la un anumit nivel, probleme care folosesc noțiuni ce nu
figurează oficial în programa acelei etape: de exemplu probleme cu derivate la OLM și la OJM, deși, nefiind în
programă, ele nu pot apărea efectiv acolo. Le-am păstrat ca exercițiu de tipul „și dacă ar pica totuși la OLM
sau la OJM'': ideea de bază este accesibilă și la acel nivel, iar antrenamentul pe ele nu poate strica.

\medskip
\noindent \textbf{Despre redactare.} Am încercat să redactez așa cum se cere la olimpiadă: enunț și demonstrez
fiecare lemă pe care o folosesc, împart soluțiile în pași numiți explicit, verific ipotezele înainte de a aplica
o teoremă și marchez clar unde se încheie fiecare argument. Redactarea nu este un moft: la concurs, ea face
diferența dintre o idee bună și punctajul maxim.

\medskip
\noindent \textbf{Despre probleme.} Aproape toate sunt inventate de mine. Excepțiile (problemele clasice și cele preluate de la concursuri) sunt indicate explicit. Am încercat, de asemenea, să includ cel puțin o problemă pentru fiecare rezultat enunțat în partea de teorie: o lemă care nu „lucrează'' nicăieri nu merită ținută minte.

\vspace{1cm}
\begin{flushright}
{\itshape Vlad-Ștefan Constantinescu}\\
{\small\textcolor{grey}{București, 2026}}
\end{flushright}
\clearpage
